% Part: first-order-logic
% Chapter: lindstrom
% Section: lindstrom-proof

\documentclass[../../../include/open-logic-section]{subfiles}

\begin{document}

\olfileid[tr]{mod}{lin}{prf}

\olsection{Lindstr\"om Teoremi}

\begin{lem}
\ollabel{lem:lindstrom}
$\Lang{L}$ sonlu olmak üzere $!E\in L(\Lang{L})$ olsun. Ayrıca öyle bir
$n\in\Nat$ bulunduğunu varsayalım ki !!{structure}s sınıfından herhangi
$\Struct{M}$ ve $\Struct{N}$ için, $\Struct{M}\equiv_n\Struct{N}$ ve
$\Struct{M}\models_L !E$ ise $\Struct{N}\models_L !E$ de olsun. O zaman
$!E$ bir birinci dereceden !!{sentence} ile eşdeğerdir; yani
$\Mod(L){!E}=\Mod(L){!D}$ olacak bir birinci dereceden $!D$ vardır.
\end{lem}

\begin{proof}
$n$, !!{structure}s sınıfından $n$-eşdeğer herhangi $\Struct{M}$ ile
$\Struct{N}$ yapısının $!E$'ye aynı doğruluk değerini verdiği bir sayı olsun.
\olref[bas][pis]{prop:qr-finite} sonucunu anımsayalım: Niceleyici derecesi
$n$'den büyük olmayan birinci dereceden !!{sentence}s, sonlu bir
!!{language} içinde mantıksal eşdeğerlik dışında yalnızca sonlu
sayıdadır. Şimdi sabit bir !!{structure}~$\Struct{M}$ için
$!D_{\Struct{M}}$, $\Struct{M}$ içinde doğru olan ve
$\QuantRank{!E}\leq n$ koşulunu sağlayan bütün birinci dereceden
!!{sentence}s~$!E$ öğelerinin birleşimi olsun. Bu birleşim sonludur ve
$\Struct{N}\models !D_{\Struct{M}}$ ancak ve ancak
$\Struct{N}\equiv_n\Struct{M}$ olduğunda geçerlidir. Ardından
\[
  !D=\textstyle\bigvee
  \Setabs{!D_{\Struct{M}}}{\Struct{M}\models_L !E}
\]
olsun. Bu ayrışım da mantıksal eşdeğerlik dışında sonludur.

$\Mod(L){!E}=\Mod(L){!D}$ olduğu artık görülür. Gerçekten,
$\Struct{N}\models !D$ ise bazı $\Struct{M}\models_L !E$ için
$\Struct{N}\models !D_{\Struct{M}}$ olur; dolayısıyla lemmanın varsayımı
uyarınca $\Struct{N}\models_L !E$ olur. Tersine,
$\Struct{N}\models_L !E$ ise $!D_{\Struct{N}}$, $!D$ ayrışımının bir
bileşenidir ve $\Struct{N}\models !D_{\Struct{N}}$ olduğundan
$\Struct{N}\models !D$ de olur.
\end{proof}

\begin{thm}[Lindstr\"om Teoremi]
  \ollabel{thm:lindstrom}
  $\tuple{L,\models_L}$ Tıkızlık ve L\"owenheim--Skolem Özelliklerine
  sahip olsun. O zaman
  $\tuple{L,\models_L}\leq\tuple{F,\models}$ olur; dolayısıyla
  $\tuple{L,\models_L}$ birinci dereceden mantığa eşdeğerdir.
\end{thm}

\begin{proof}
\olref{lem:lindstrom} uyarınca, sonlu $\Lang{L}$ için her
$!E\in L(\Lang{L})$ verildiğinde öyle bir $n\in\Nat$ bulunduğunu göstermek
yeterlidir ki !!{structure}s sınıfından herhangi $\Struct{M}$ ile
$\Struct{N}$ için
$\Struct{M}\equiv_n\Struct{N}$ olması, bu iki yapının $!E$ üzerinde
uzlaşmasını gerektirsin. O zaman $!E$ bir birinci dereceden !!{sentence}
ile eşdeğer olur ve buradan
$\tuple{L,\models_L}\leq\tuple{F,\models}$ sonucu çıkar. Sonlu ve bütünüyle
bağıntısal bir !!{language} içinde çalıştığımızdan,
\olref[bas][pis]{thm:b-n-f} uyarınca
$\Struct{M}\equiv_n\Struct{N}$ ifadesini buna karşılık gelen cebirsel
$I_n(\emptyset,\emptyset)$ ifadesiyle değiştirebiliriz.

Bir $!E$ verilsin. Çelişkiye ulaşmak üzere, her $n$ için
$I_n(\emptyset,\emptyset)$ koşulunu sağlayan, fakat örneğin
$\Struct{M}_n\models_L !E$ ve $\Struct{N}_n\not\models_L !E$ olan
!!{structure}s $\Struct{M}_n$ ile $\Struct{N}_n$ bulunduğunu varsayalım.
İzomorfizm Özelliği sayesinde bütün $\Struct{M}_n$ yapılarının dildeki
sabitleri aynı nesneler olarak yorumladığını varsayabiliriz. Üstelik dilde
yalnızca sonlu sayıda atomik !!{sentence}s bulunduğundan, bütün bu yapıların
aynı atomik !!{sentence}s öğelerini sağladığını da varsayabiliriz; gerekirse
$\Struct{M}_n$ dizisinin bir alt dizisini alırız. $\Struct{M}$, bütün
$\Struct{M}_n$ yapılarının birleşimi, yani her $\Struct{M}_n$ yapısını alt
yapı olarak içeren tek en küçük !!{structure} olsun. Şimdi
\olref[lsp]{thm:abstract-p-isom} ispatındaki gibi, $\Struct{M}$ yapısını
!!{domain} $\Domain{M}\cup\Domain{M}^{<\omega}$ olacak ve genişletilmiş
!!{language} içinde birleştirme yüklemleri $P$ ile $Q$'yu içerecek biçimde
$\Struct{M}^*$ yapısına genişletelim.

$\Struct{N}_n$, $\Struct{N}$ ve $\Struct{N}^*$ benzer biçimde tanımlansın.
Şimdi bir $\Struct{K}$ !!{structure} kuralım. Bu yapının !!{domain},
$\Struct{M}^*$ ile $\Struct{N}^*$ yapılarının !!{domain}s öğelerinin yanı
sıra doğal sıralamaları $\leq$ ile birlikte doğal sayıları $\Nat$ içersin.
!!{language} dili, $\Domain{M}$, $\Domain{N}$,
$\Domain{M}^{<\omega}$ ve $\Domain{N}^{<\omega}$ !!{domain}s öğelerini
temsil eden ek yüklemler ile $\Struct{M}_n$ ve $\Struct{N}_n$ evrenlerini
aşağıdaki anlamda kodlayan yüklemler içersin:
\begin{align*}
  \Domain{M_n} & = \Setabs{a\in\Domain{M}}{R(a,n)}; &
  \Domain{N_n} & = \Setabs{a\in\Domain{N}}{S(a,n)}; \\
  \Domain{M}^{<\omega}_n & =
    \Setabs{a\in\Domain{M}^{<\omega}}{R(a,n)}; &
  \Domain{N}^{<\omega}_n & =
    \Setabs{a\in\Domain{N}^{<\omega}}{S(a,n)}.
\end{align*}
!!{structure}~$\Struct{K}$ ayrıca, $J(n,\mathbf a,\mathbf b)$ ancak ve
ancak $I_n(\mathbf a,\mathbf b)$ olduğunda geçerli olacak bir üçlü $J$
bağıntısına sahip olsun.

Şimdi öyle !!a{sentence}~$!D$ vardır ki $R$, $S$, $J$ vb. simgelerle
genişletilmiş !!{language} içinde $\leq$ bağıntısının ilk elemanı bulunan fakat
son elemanı bulunmayan ayrık bir tam sıralama olduğunu;
$\Struct{M}_n\models_L !E$ ve $\Struct{N}_n\not\models_L !E$ olduğunu; ayrıca
sıralamadaki her $n$ için $J(n,\mathbf a,\mathbf b)$ ancak ve ancak
$I_n(\mathbf a,\mathbf b)$ olduğunu söyler.

Tıkızlık Özelliğini kullanarak $!D$'nin, sıralamasında standart olmayan bir
$n^*$ elemanı bulunan bir $\Struct{K}^*$ modelini elde edebiliriz. Özellikle
$\Struct{K}^*$, $\Struct{M}_{n^*}\models_L !E$ ve
$\Struct{N}_{n^*}\not\models_L !E$ olacak biçimde alt!!{structure}s
$\Struct{M}_{n^*}$ ile $\Struct{N}_{n^*}$ yapısını içerir. Şimdi
$\Domain{M_{n^*}}$ ile $\Domain{N_{n^*}}$ evrenlerinden gelen $k$-lilerin
çiftlerinden oluşan bir $\mathcal I$ kümesini şöyle tanımlayabiliriz:
$\tuple{\mathbf a,\mathbf b}\in\mathcal I$ ancak ve ancak
$J(n^*-k,\mathbf a,\mathbf b)$ olsun; burada $k$, $\mathbf a$ ile
$\mathbf b$ dizilerinin uzunluğudur. $n^*$ standart olmadığından, her
standart $k$ için $n^*-k>0$ olur ve $\mathcal I$ kümesi
$\Struct{M}_{n^*}\simeq_p\Struct{N}_{n^*}$ olduğunu gösterir. Oysa
\olref[lsp]{thm:abstract-p-isom} uyarınca $\Struct{M}_{n^*}$ ile
$\Struct{N}_{n^*}$ yapıları $L$-eşdeğerdir. Bu bir çelişkidir.
\end{proof}

\end{document}
